\section{技術スタック}

\subsection{選定方針}

\begin{itemize}[leftmargin=*,itemsep=2pt]
  \item \textbf{型安全と開発体験}: TypeScript + Vite による高速な
        HMR と厳格な型チェック。
  \item \textbf{標準準拠}: Canvas 2D / FontFace / LocalStorage など
        W3C 標準 API のみを使用し、フレームワーク依存を最小化。
  \item \textbf{UI 構築速度}: Tailwind CSS + shadcn/ui（Radix UI ベース）
        でスタイルと挙動を分離しつつ、アクセシビリティを確保。
  \item \textbf{ゼロ運用コスト}: GitHub Pages による静的配信。
        Actions で自動デプロイ、サーバは持たない。
\end{itemize}

\subsection{使用ライブラリ・ツール}

\begin{tabularx}{\textwidth}{l l X}
\toprule
\textbf{分類} & \textbf{名称} & \textbf{用途} \\
\midrule
言語       & TypeScript 5     & 型安全な JavaScript。\\
ランタイム & Node.js 20+      & 開発・ビルド環境。\\
ビルド     & Vite 6           & 開発サーバ + 本番バンドル。\\
UI 基盤    & React 18         & 宣言的 UI、Hooks ベース状態管理。\\
スタイル   & Tailwind CSS 3   & ユーティリティファースト CSS。\\
UI 部品    & shadcn/ui        & Radix UI + Tailwind プリミティブ集。\\
UI プリミティブ & Radix UI    & アクセシブルな Tabs/Select/Slider 等。\\
アイコン   & lucide-react     & SVG アイコン。\\
CVA        & class-variance-authority & バリアント管理。\\
\bottomrule
\end{tabularx}

\subsection{ブラウザ API}

\begin{tabularx}{\textwidth}{l X}
\toprule
\textbf{API} & \textbf{役割} \\
\midrule
Canvas 2D (\texttt{getContext("2d")})     & フォント描画とピクセル抽出。\\
\texttt{getImageData}                     & RGBA 配列取得、アルファ閾値 2 値化の基礎。\\
\texttt{FontFace}                         & カスタム TTF/OTF/WOFF の動的登録。\\
\texttt{URL} / \texttt{URLSearchParams}   & 共有リンクのクエリパラメータ読み書き。\\
\texttt{LocalStorage}                     & 設定の永続化。\\
\texttt{navigator.clipboard}              & 出力結果のコピー。\\
\texttt{matchMedia("prefers-color-scheme")} & OS テーマ検出。\\
\bottomrule
\end{tabularx}

\subsection{同梱フォント}

\begin{tabularx}{\textwidth}{l l X}
\toprule
\textbf{フォント} & \textbf{推奨サイズ} & \textbf{ライセンス} \\
\midrule
DotGothic16       & 16$\times$16 & SIL OFL 1.1 \\
Misaki Gothic 2nd & 8$\times$8   & 自由配布（商用含む）\\
システムモノスペース & 任意       & ブラウザ標準フォント依存 \\
\bottomrule
\end{tabularx}

\subsection{ビルド設定}

\begin{itemize}[leftmargin=*,itemsep=2pt]
  \item \texttt{vite.config.ts}: \texttt{base} を本番時のみ
        \texttt{/font\_to\_bin/} に切り替え。
  \item \texttt{tsconfig.json}: \texttt{strict: true}, \texttt{noUnusedLocals},
        \texttt{noUnusedParameters} を有効化。
  \item \texttt{tailwind.config.js}: \texttt{darkMode: "class"}、
        カスタムカラートークンを定義。
  \item PWA manifest (\texttt{public/manifest.webmanifest}):
        アイコン・テーマカラー・表示モードを定義（SW は持たない）。
\end{itemize}

\newpage
